\clearpage
\VerdierPhysicalPageAnchor{0188}%
\begingroup\small\setlength{\parskip}{0pt}\setlength{\emergencystretch}{1.5em}
\setlength{\abovedisplayskip}{4pt plus 2pt minus 2pt}%
\setlength{\belowdisplayskip}{4pt plus 2pt minus 2pt}%

\VerdierSemanticAnchor{ch3d-2-1-3-properties-cont}
初等射である。初等射は基底変換で安定である。次の複体の射の列を考える：

\VerdierSemanticAnchor{ch3d-2-1-3-bi-infinite-chain}
\[
\cdots\longrightarrow X_n\xrightarrow{u_n}X_{n-1}
\xrightarrow{u_{n-1}}X_{n-2}\longrightarrow\cdots
\]
この列は両方向に無限であり、各整数 $n$ に対して射 $u_n$ は高さ $-n$ の
初等射であるとする。この列は $\operatorname{comp}(\mathcal A)$ において
射影極限と帰納極限をもつ（任意に固定した次数の成分に誘導される列は、
両方向で定常である）。

\VerdierSemanticAnchor{ch3d-2-1-4}
\medskip\noindent\textbf{補題 2.1.4.}\quad
\textit{$m\colon X\longrightarrow Y$ を
$\operatorname{comp}(\mathcal A)$ の擬同型とする。射の列}
\[
u_n\colon X_n\longrightarrow X_{n-1},\qquad n\in\mathbb Z,
\]
\textit{で、各整数 $n$ に対して $u_n$ が高さ $-n$ の初等射となるもの、および
射}

\VerdierSemanticAnchor{ch3d-2-1-4-limits}
\[
\varinjlim_n X_n\xleftarrow{\ \sim\ }Y,
\qquad
\varprojlim_n X_n\xleftarrow{\ v\ }X,
\quad v\text{ はホモトピー同値（第 I 章，2.5.11）}
\]
\textit{であって、図式}

\VerdierSemanticAnchor{ch3d-2-1-4-square}
\begin{center}
\begin{tikzcd}[column sep=5.6em,row sep=3em]
X \arrow[r,"v"] \arrow[d,"m"']
  & \displaystyle\varprojlim_n X_n \arrow[d] \\
Y \arrow[r,"{\sim}"'] & \displaystyle\varinjlim_n X_n
\end{tikzcd}
\end{center}
\textit{がホモトピーを除いて可換となるものが存在する。}

\VerdierSemanticAnchor{ch3d-2-1-4-proof}
まず、射 $m\colon X\longrightarrow Y$ を
$\operatorname{comp}(\mathcal A)$ の射とすると、エピ射
$m'\colon X'\longrightarrow Y$ と、ホモトピーを除いて可換な図式

\VerdierSemanticAnchor{ch3d-2-1-4-epimorphism-diagram}
\begin{center}
\begin{tikzcd}[column sep=3.5em,row sep=3.2em]
X \arrow[rr,"v"] \arrow[dr,"m"'] && X' \arrow[dl,"{m'}"] \\
& Y &
\end{tikzcd}
\end{center}
が存在する。
\endgroup
